% Quick preview/compile notes:
% - Local (Windows): install MiKTeX or TeX Live. In CMD:
%     cd d:\CODE_AIC_2025
%     pdflatex -interaction=nonstopmode -halt-on-error test.tex
%   Run pdflatex twice if cross-references need updating.
% - Overleaf: create project, upload this file, press "Recompile".
% - Editor with live preview: use TeXstudio or VS Code + LaTeX Workshop.
% - If a package (e.g., chemfig) is missing, install it via MiKTeX/TeX Live manager.
% Editor directives (helpful for TeXstudio/VSCode):
% !TeX program = pdflatex
% !TeX encoding = UTF-8
\documentclass{article}
\usepackage[utf8]{inputenc}  % Để hỗ trợ tiếng Việt
\usepackage{chemfig}
\usepackage{color}

\begin{document}

\begin{figure}[htbp]
\centering
\schemestart
  % Phần chất phản ứng
  \subscheme{
    \chemfig{\textcolor{red}{HO^-}}
    \arrow{0}[0,0.5]  % Khoảng cách
    \chemfig{CH_3-@{c}C(-[6]H)-@{cl}Cl}
    \chemmove{
      \draw[black,-stealth,shorten <=2pt,shorten >=2pt] (HO.east) .. controls +(0:1cm) and +(180:1cm) .. (c.west);
      \draw[black,-stealth,shorten <=2pt,shorten >=2pt] (c.east) .. controls +(0:0.5cm) and +(180:0.5cm) .. (cl.west);
      \node[above left=0.1cm] at (c) {\textcolor{black}{\delta^+}};
      \node[above right=0.1cm] at (cl) {\textcolor{black}{\delta^-}};
    }
  }
  \arrow{->[ ][NaOH]}[0,1.5]  % Mũi tên với NaOH trên

  % Phần trạng thái chuyển tiếp (với ngoặc vuông)
  \subscheme{
    \chemleft[ \chemfig{\textcolor{red}{HO^{\delta^-}} -[,,,,dotted] @{c}C^{\delta^+} (-[2]CH_3)(-[6]H) -[,,,,dotted] {\textcolor{blue}{Cl^{\delta^-}}}} \chemright]
    \chemmove{
      \node[below=0.2cm] at (c) {H \quad H};  % Thêm hai H bên dưới nếu hình gốc có
    }
  }
  \arrow{->}[0,1.5]

  % Phần sản phẩm phản ứng
  \subscheme{
    \chemfig{\textcolor{red}{HO} -C(-[2]CH_3)(-[6]H)-H} + \textcolor{blue}{Cl^-}
  }

\schemestop

% Thêm nhãn tiếng Việt bên dưới từng phần
\chemnameinit{}  % Reset chemname
\begin{center}
  \begin{tabular}{ccc}
    chất phản ứng & trạng thái chuyển tiếp & sản phẩm phản ứng \\
  \end{tabular}
\end{center}

\caption{Cơ chế phản ứng SN2.}
\label{fig:co_che_sn2}
\end{figure}

\end{document}